% GENERATED FILE. Edit canonical lessons, then run npm run generate:books.
% canonical-source-hash: fnv1a64:121249aa5c441585
% canonical-lessons: IT-W262-four-lines

\chapter{Writing --- Four Lines, And The Full Stop That Answers}
\label{ch:it-four-lines}
\hlchaptermodality{262}

\begin{chapteropening}
\textbf{By the end of this chapter:} \emph{I can write four connected lines to a named reader, joining with e and ma and giving a reason with perché, ending that line with a full stop rather than a question mark.}

Complete IT-W262-four-lines: four lines to one reader, at least one joined with e or ma and one with perché, then check the perché line ends with a full stop and its accent rises.
\end{chapteropening}

\section[(four lines)]{Four lines, and the full stop that answers}

\label{lesson:IT-W262-four-lines}



\begin{quote}
The last lesson asked for writing \textbf{finished on time}. This one asks for writing that \textbf{holds together}.

Pick one person who will read it.
\end{quote}

\subsection*{What to know first}
Connected writing is not \textbf{longer} writing. Four short lines can be connected; one long line can fail to be. What makes it connected is that a sentence stops standing on its own.

You hold three joiners. \textbf{e} adds. \textbf{ma} sets against. \textbf{perché} gives the reason. After \textbf{perché} the words keep the order they would have on their own: \emph{perché non mi piace}. Nothing moves.

And \textbf{perché} is one word for two jobs: \emph{why} and \emph{because}. Spoken, the voice tells them apart. On paper only the end of the line does: a \textbf{question mark} asks, a \textbf{full stop} answers.

Its accent leans the other way from the one on \textbf{è}: \textbf{é} rises as it goes right, \textbf{è} falls as it goes right. One line can hold both, so write each one on purpose.

\subsection*{Writing — connected composition}
Write \textbf{four lines} to the person you picked, on paper.

At least one must use \textbf{e} or \textbf{ma}. At least one must give a reason with \textbf{perché}, and end with a full stop.

Use only words you already write. This lesson adds no vocabulary; the whole exercise is the joining.

\begin{tcolorbox}[breakable,colback=teal!4,colframe=teal!35!black,title={Before you move on}]
Cover your first line and read the second alone. Does it still make sense without the first? If it does, you wrote a \textbf{list}, which is where you were one lesson ago.

Then find your \textbf{perché} line. A full stop at its end, not a question mark? An accent that rises? If either is wrong, you wrote a question, or a different mark, and the reader will meet it before they meet your reason.
\end{tcolorbox}
